\begin{thebibliography}{99}\small
\bibitem{jermyn23} Jermyn, A. S., et al. (2023).
Modules for Experiments in Stellar Astrophysics (MESA): Time-Dependent
Convection, Energy Conservation, Automatic Differentiation, and Infrastructure.
\emph{ApJS}, 265, 15. \url{https://arxiv.org/abs/2208.03651}.
\bibitem{mesa_atm} MESA documentation. Table atmospheres.
\url{https://docs.mesastar.org/en/latest/atm/table.html}.
\bibitem{rohrmann11} Rohrmann, R. D., Althaus, L. G., \& Kepler, S. O. (2011).
Lyman $\alpha$ wing absorption in cool white dwarf stars.
\emph{MNRAS}, 411, 781--791.
\href{https://doi.org/10.1111/j.1365-2966.2010.17716.x}{doi:10.1111/j.1365-2966.2010.17716.x}.
\bibitem{colgan16} Colgan, J., et al. (2016).
A New Generation of Los Alamos Opacity Tables. \textit{ApJ}, 817, 116.
\href{https://doi.org/10.3847/0004-637X/817/2/116}{doi:10.3847/0004-637X/817/2/116}.
\bibitem{tops} Los Alamos National Laboratory. TOPS opacity service,
ATOMIC/New elemental library and mixture queries.
\href{https://aphysics2.lanl.gov/opacity/lanl/}{aphysics2.lanl.gov/opacity/lanl/}.
\bibitem{davis24} Davis, Y. T., et al. (2024).
The EBLM Project XII. An eccentric, long-period eclipsing binary with a companion
near the hydrogen-burning limit.
\href{https://arxiv.org/pdf/2312.09156v2}{arXiv:2312.09156v2}, Table 2 and abstract.
\bibitem{agol21} Agol, E., et al. (2021).
Refining the Transit-timing and Photometric Analysis of TRAPPIST-1:
Masses, Radii, Densities, Dynamics, and Ephemerides.
\href{https://arxiv.org/pdf/2010.01074}{arXiv:2010.01074}, Table 7.
\bibitem{ribas17} Ribas, I., et al. (2017).
The full spectral radiative properties of Proxima Centauri.
\href{https://arxiv.org/pdf/1704.08449}{arXiv:1704.08449}, Table 11.
\bibitem{lba97} Laughlin, G., Bodenheimer, P., \& Adams, F. C. (1997).
The End of the Main Sequence. \emph{ApJ}, 482, 420.
\href{https://doi.org/10.1086/304125}{doi:10.1086/304125}.
\bibitem{chabrier97} Chabrier, G., \& Baraffe, I. (1997).
Structure and evolution of low-mass stars. \emph{A\&A}, 327, 1039.
\href{https://arxiv.org/abs/astro-ph/9704118}{arXiv:astro-ph/9704118}, Table 2.
\bibitem{adams97} Adams, F. C., \& Laughlin, G. (1997).
A dying universe: the long-term fate and evolution of astrophysical objects.
\emph{Rev. Mod. Phys.}, 69, 337.
\href{https://doi.org/10.1103/RevModPhys.69.337}{doi:10.1103/RevModPhys.69.337}.
\bibitem{rathore05} Rathore, Y., Blandford, R. D., \& Broderick, A. E. (2005).
Resonant excitation of white dwarf oscillations in compact object binaries---I.
The no back reaction approximation.
\emph{MNRAS}, 357, 834.
\href{https://academic.oup.com/mnras/article/357/3/834/1077414}{Published calculation}.
\bibitem{hrw94} Haft, M., Raffelt, G., \& Weiss, A. (1994).
Standard and nonstandard plasma neutrino emission revisited.
\href{https://arxiv.org/abs/astro-ph/9309014}{arXiv:astro-ph/9309014}.
\bibitem{chabrier00} Chabrier, G., Baraffe, I., Allard, F., \& Hauschildt, P. (2000).
Evolutionary models for very low-mass stars and brown dwarfs with dusty atmospheres.
\emph{ApJ}, 542, 464.
\href{https://arxiv.org/abs/astro-ph/0005557}{arXiv:astro-ph/0005557}.
\bibitem{baiko22} Baiko, D. A., \& Chugunov, A. I. (2022).
\emph{Ab initio} thermodynamics of one-component plasma for astrophysics of
white dwarfs and neutron stars.
\href{https://arxiv.org/abs/2112.04822}{arXiv:2112.04822}.
\bibitem{blouin18} Blouin, S., Dufour, P., \& Allard, N. F. (2018).
A new generation of cool white dwarf atmosphere models. I. Theoretical framework
and applications to DZ stars.
\href{https://arxiv.org/abs/1807.06616}{arXiv:1807.06616}.
\bibitem{beznogov16} Beznogov, M. V., Potekhin, A. Y., \& Yakovlev, D. G. (2016).
Diffusive heat blanketing envelopes of neutron stars.
\emph{MNRAS}, 459, 1569.
\href{https://doi.org/10.1093/mnras/stw751}{doi:10.1093/mnras/stw751}.
\bibitem{thoul94} Thoul, A. A., Bahcall, J. N., \& Loeb, A. (1994).
Element diffusion in the solar interior. ApJ, 421, 828.
\href{https://arxiv.org/abs/astro-ph/9304005}{arXiv:astro-ph/9304005}.
\bibitem{denissenkov13} Denissenkov, P. A., Herwig, F., Bildsten, L.,
\& Paxton, B. (2013). MESA Models of Classical Nova Outbursts: The Multicycle
Evolution and Effects of Convective Boundary Mixing. \emph{ApJ}, 762, 8.
\href{https://doi.org/10.1088/0004-637X/762/1/8}{Published article}.
\bibitem{bauer19} Bauer, E. B., \& Bildsten, L. (2019).
Polluted White Dwarfs: Mixing Regions and Diffusion Timescales.
\emph{ApJ}, 872, 96. \href{https://arxiv.org/abs/1812.09602}{arXiv:1812.09602}.
\bibitem{charbonnel07} Charbonnel, C., \& Zahn, J.-P. (2007).
Thermohaline mixing: A physical mechanism governing the photospheric
composition of low-mass giants. \emph{A\&A}, 467, L15.
\href{https://arxiv.org/abs/astro-ph/0703302}{arXiv:astro-ph/0703302}.
\bibitem{paxton11} Paxton, B., et al. (2011).
Modules for Experiments in Stellar Astrophysics (MESA).
\emph{ApJS}, 192, 3.
\href{https://doi.org/10.1088/0067-0049/192/1/3}{Published methods}.
\bibitem{paxton18} Paxton, B., et al. (2018).
Modules for Experiments in Stellar Astrophysics (MESA): Convective Boundaries,
Element Diffusion, and Massive Star Explosions.
\href{https://arxiv.org/abs/1710.08424}{arXiv:1710.08424}, section 3.2.2.
\bibitem{adelberger11} Adelberger, E. G., et al. (2011).
Solar fusion cross sections. II. The pp chain and CNO cycles.
Reviews of Modern Physics, 83, 195.
\bibitem{acharya25} Acharya, B., et al. (2025).
Solar fusion III: New data and theory for hydrogen-burning stars.
\emph{Rev. Mod. Phys.}, 97, 035002.
\href{https://doi.org/10.1103/8lm7-gs18}{doi:10.1103/8lm7-gs18}.
\bibitem{baraffe15} Baraffe, I., Homeier, D., Allard, F., \& Chabrier, G. (2015).
New evolutionary models for pre-main sequence and main sequence low-mass stars
down to the hydrogen-burning limit.
\href{https://arxiv.org/abs/1503.04107}{arXiv:1503.04107}.
\bibitem{takenaka20} Takenaka, A., et al. (Super-Kamiokande Collaboration, 2020).
Search for proton decay via $p\to e^+\pi^0$ and $p\to\mu^+\pi^0$
with an enlarged fiducial volume in Super-Kamiokande I--IV.
\href{https://arxiv.org/abs/2010.16098}{arXiv:2010.16098}.
\bibitem{fukugita86} Fukugita, M., \& Yanagida, T. (1986).
Baryogenesis without grand unification. \emph{Phys. Lett. B}, 174, 45.
\href{https://doi.org/10.1016/0370-2693(86)91126-3}{doi:10.1016/0370-2693(86)91126-3}.
\bibitem{adams98} Adams, F. C., Laughlin, G., Mbonye, M., \& Perry, M. J. (1998).
The gravitational demise of cold degenerate stars. \emph{Phys. Rev. D}, 58, 083003.
\href{https://arxiv.org/abs/astro-ph/9808250}{arXiv:astro-ph/9808250}.
\bibitem{adams01} Adams, F. C., Kane, G. L., Mbonye, M., \& Perry, M. J. (2001).
Proton decay, black holes, and large extra dimensions.
\href{https://arxiv.org/abs/hep-ph/0009154}{arXiv:hep-ph/0009154}.
\bibitem{bell24} Bell, N. F., et al. (2024).
Heavy Dark Matter in White Dwarfs: Multiple-Scattering Capture and Thermalization.
\href{https://arxiv.org/abs/2404.16272}{arXiv:2404.16272}.
\bibitem{gasques05} Gasques, L. R., et al. (2005).
Nuclear fusion in dense matter: Reaction rate and carbon burning.
\emph{Phys. Rev. C}, 72, 025806.
\href{https://arxiv.org/abs/astro-ph/0506386}{arXiv:astro-ph/0506386}.
\bibitem{lederer09} Lederer, M. T., \& Aringer, B. (2009).
Low temperature Rosseland opacities with varied abundances of carbon and nitrogen.
\href{https://doi.org/10.1051/0004-6361:200810576}{Published article};
\href{https://arxiv.org/abs/0810.5672}{arXiv:0810.5672}.
\bibitem{potekhin12} Potekhin, A. Y., \& Chabrier, G. (2012).
Thermonuclear fusion in dense stars: Electron screening, conductive cooling,
and magnetic field effects. A\&A, 538, A115.
\href{https://arxiv.org/abs/1201.2133}{arXiv:1201.2133}.
\bibitem{chugunov09} Chugunov, A. I., \& DeWitt, H. E. (2009).
Nuclear fusion reaction rates for strongly coupled ionic mixtures.
Physical Review C, 80, 014611.
\href{https://arxiv.org/abs/0905.3844}{arXiv:0905.3844}.
\bibitem{seaton05} Seaton, M. J. (2005).
OP data on CD for mean opacities and radiative accelerations.
\href{https://doi.org/10.1111/j.1365-2966.2005.00019.x}{doi:10.1111/j.1365-2966.2005.00019.x}.
\bibitem{farag24} Farag, E., et al. (2024).
An Expanded Set of Los Alamos OPLIB Tables in MESA:
Type-1 Rosseland-mean Opacities and Solar Models.
\textit{The Astrophysical Journal}.
\href{https://doi.org/10.3847/1538-4357/ad4355}{doi:10.3847/1538-4357/ad4355}.
\end{thebibliography}
